% <<TITLE>>: <<SUBTITLE>>
% Build: bash ~/.agents/skills/research-report/scripts/build.sh <this file>  (runs xelatex twice; fonts in fonts/, OFL)
% Source of truth: <<PATH TO RESEARCH NOTE>>. Every claim here must exist in that note first.
% Citing: \src{S1} makes a clickable [S1] that jumps to the source list; each source there links to its URL.
% To add a source: add a \source{Sx}{citation}{url} line at the bottom, then cite it with \src{Sx}.
\documentclass[11pt,a4paper]{article}

\usepackage[margin=2.4cm]{geometry}
\usepackage{fontspec}
\setmainfont{Inter}[Path=fonts/, Extension=.otf, UprightFont=*-Regular, BoldFont=*-Bold,
  ItalicFont=*-Italic, BoldItalicFont=*-BoldItalic]
\usepackage{microtype}
\usepackage[table]{xcolor}
\usepackage{booktabs, tabularx, array}
\usepackage{titlesec, enumitem, fancyhdr, tcolorbox}
\usepackage{parskip}
\usepackage{xurl}
\usepackage{tikz}

% Colours. Default teal; use the project's brand colours when it has them
% (accent: a dark brand colour that reads as small text; soft: a pale tint for boxes; muted: secondary text).
\definecolor{accent}{HTML}{0F6E6E}
\definecolor{soft}{HTML}{EEF5F5}
\definecolor{muted}{HTML}{5B6770}
\definecolor{ink}{HTML}{0F1220}

\usepackage[colorlinks=true, linkcolor=accent, urlcolor=accent, citecolor=accent,
  pdftitle={<<TITLE>>: <<SUBTITLE>>}]{hyperref}

\linespread{1.12}
\setlist{itemsep=0.25em, topsep=0.3em}
\renewcommand{\arraystretch}{1.25}
\setlength{\tabcolsep}{6pt}

\titleformat{\section}{\Large\bfseries\color{accent}}{}{0pt}{}[\vspace{-0.4em}{\color{accent}\titlerule[0.6pt]}]
\titleformat{\subsection}{\large\bfseries}{}{0pt}{}
\titlespacing*{\section}{0pt}{1.6em}{0.6em}
\titlespacing*{\subsection}{0pt}{1.1em}{0.3em}

\pagestyle{fancy}
\fancyhf{}
\renewcommand{\headrulewidth}{0pt}
\fancyfoot[L]{\footnotesize\color{muted}<<SHORT TITLE>> \textbar{} draft <<YYYY-MM-DD>>}
\fancyfoot[R]{\footnotesize\color{muted}\thepage}

% Citations: clickable [S1] linking to the source list.
\newcommand{\src}[1]{\hyperlink{src:#1}{[#1]}}
% One entry in the source list: key, citation text, URL.
\newcommand{\source}[3]{\item[\hypertarget{src:#1}{\textbf{\textcolor{accent}{[#1]}}}] #2 \\ \url{#3}}
\newenvironment{sources}{\begin{description}[leftmargin=2.6em, labelwidth=2.6em, itemsep=0.5em, font=\normalfont]}{\end{description}}

% Evidence level after a claim, e.g. \lvl{vendor survey}; our own reasoning gets \inferred.
\newcommand{\lvl}[1]{\textcolor{muted}{(#1)}}
\newcommand{\inferred}{\textcolor{muted}{\emph{(inferred)}}}
\newtcolorbox{callout}[1][]{colback=soft, colframe=accent, boxrule=0pt, leftrule=3pt,
  arc=0pt, left=10pt, right=10pt, top=6pt, bottom=6pt, #1}

% Logo, top right of page 1. Empty by default. Fill it only from the project's design source
% (an SVG file or brand config), never from a screenshot. scripts/svg2tikz.py turns SVG paths into
% the \fill lines; see SKILL.md "Logo". Example with a mark plus a name in a brand font:
%   \newfontfamily\brandtext{BrandFont-SemiBold}[Path=fonts/, Extension=.ttf]
%   \newcommand{\logomark}[1]{\begin{tikzpicture}[x=1pt, y=1pt, scale=#1/24, baseline={#1pt*4.2/24}]
%     \fill[ink] ... ;   % output of svg2tikz.py
%   \end{tikzpicture}}
%   \newcommand{\reportlogo}{\mbox{\logomark{15.4}\hspace{6.2pt}{\brandtext\fontsize{14}{14}\selectfont\color{ink}Name}}}
\newcommand{\reportlogo}{}

\begin{document}

% ---------------------------------------------------------------- title
{\color{accent}\bfseries\small DRAFT FOR REVIEW \quad\textbar\quad <<DATE IN CAPITALS, e.g. 25 NOVEMBER 2026>>\hfill\normalcolor\reportlogo\par}
\vspace{0.4em}
{\fontsize{26}{30}\selectfont\bfseries <<TITLE>>\par}
{\fontsize{15}{19}\selectfont\color{muted} <<SUBTITLE>>\par}
\vspace{0.6em}
{\small\color{muted} <<WHO IT IS WRITTEN FOR, ONE SENTENCE>>. Every claim carries a clickable source reference and an
evidence level. Our own reasoning is marked \emph{(inferred)}.\par}

% ---------------------------------------------------------------- open questions
% Only the note's open questions for review. Drop the box if the note has none.
\begin{callout}
\textbf{Open questions for review} (none blocked this draft)
\begin{enumerate}[leftmargin=1.6em, itemsep=0.15em]
  \item \textbf{<<Short label.>>} <<One or two sentences.>>
\end{enumerate}
\end{callout}

% ---------------------------------------------------------------- answer first
\section{Answer first}
\begin{itemize}
  \item \textbf{<<Lead.>>} <<The note's answer-first bullets, with numbers and \src{S1}.>>
  \item \textbf{Confidence: <<high|medium|low>>.} <<One-line reason, from the note.>>
\end{itemize}

% ---------------------------------------------------------------- scope
\section{Question and scope}
<<The note's sharpened question: question, type, scope, out of scope. If the question was reframed, say how and why.>>

% ---------------------------------------------------------------- evidence
% One numbered section per part of the note's evidence, in the note's order and with its headings.
\section{1. <<Evidence heading from the note>>}
\begin{itemize}
  \item \textbf{<<Claim.>>} <<Detail>> \src{S1}. \lvl{<<evidence level>>}
\end{itemize}

% A table only for figures the note already lists side by side. No new columns, no mapping rows.
% \begin{center}
% \begin{tabularx}{0.8\textwidth}{>{\raggedright\arraybackslash}X r}
% \toprule
% Finding & Value \\
% \midrule
% <<row>> & <<value>> \\
% \bottomrule
% \end{tabularx}
% \end{center}
% {\footnotesize\color{muted} Source: \src{S1}.\par}

% ---------------------------------------------------------------- limits
\section{Alternative explanations and limits}
\begin{itemize}
  \item <<From the note.>>
\end{itemize}

\section{Open questions}
\begin{itemize}
  \item <<From the note.>>
\end{itemize}

% ---------------------------------------------------------------- sources
\section{Sources}
\begin{sources}
\source{S1}{<<Author. ``Title''. Publisher, date. Primary|Secondary; evidence level.>>}{<<https://...>>}
\end{sources}

\textbf{Not opened:} <<sources the note lists as not read, or drop this line>>.

\end{document}
